\documentclass[11pt]{article}

% Change "review" to "final" to generate the final (sometimes called camera-ready) version.
% Change to "preprint" to generate a non-anonymous version with page numbers.
\usepackage[]{acl}

% Standard package includes
\usepackage{times}
\usepackage{latexsym}

% For proper rendering and hyphenation of words containing Latin characters (including in bib files)
\usepackage[T1]{fontenc}
% For Vietnamese characters
% \usepackage[T5]{fontenc}
% See https://www.latex-project.org/help/documentation/encguide.pdf for other character sets

% This assumes your files are encoded as UTF8
\usepackage[utf8]{inputenc}

% This is not strictly necessary, and may be commented out,
% but it will improve the layout of the manuscript,
% and will typically save some space.
\usepackage{microtype}

% This is also not strictly necessary, and may be commented out.
% However, it will improve the aesthetics of text in
% the typewriter font.
\usepackage{inconsolata}

%Including images in your LaTeX document requires adding
%additional package(s)
\usepackage{graphicx}

\usepackage{amsmath}
% If the title and author information does not fit in the area allocated, uncomment the following
%
%\setlength\titlebox{<dim>}
%
% and set <dim> to something 5cm or larger.

\title{IFT6289-H26 - Project Proposal\\ Lexical Functions to Improve Lexical
  Competence}

% Author information can be set in various styles:
% For several authors from the same institution:
% \author{Author 1 \and ... \and Author n \\
%         Address line \\ ... \\ Address line}
% if the names do not fit well on one line use
%         Author 1 \\ {\bf Author 2} \\ ... \\ {\bf Author n} \\
% For authors from different institutions:
% \author{Author 1 \\ Address line \\  ... \\ Address line
%         \And  ... \And
%         Author n \\ Address line \\ ... \\ Address line}
% To start a separate ``row'' of authors use \AND, as in
% \author{Author 1 \\ Address line \\  ... \\ Address line
%         \AND
%         Author 2 \\ Address line \\ ... \\ Address line \And
%         Author 3 \\ Address line \\ ... \\ Address line}

\author{
    \textbf{Nathalie Ann\textsuperscript{1}},
    \textbf{Gabriel Beaudoin\textsuperscript{2}},
    \textbf{Noah Tremblay Taillon\textsuperscript{3}}
    \\
    \textsuperscript{1} @umontreal.ca,
    \textsuperscript{2} gabriel.beaudoin.3@umontreal.ca,\\
    \textsuperscript{3} noah.tremblay.taillon@umontreal.ca
    }

\begin{document}
\maketitle
\begin{abstract}
Large language models exhibit strong performance across many natural language
processing tasks, yet their understanding of fine-grained lexical relations
remains limited compared to that of human native speakers. Recent work has used 
the Meaning-Text Theory (MTT) to evaluate models' lexical competence through
lexical functions, but has focused primarily on evaluation rather than
on model improvement. In this project, we propose a structured dataset of French
lexical units annotated with lexical functions derived from the French Lexical
Network. We use this dataset to fine-tune and evaluate three language models,
assessing whether explicit training on a structured description of natural
language improves models' lexical understanding and their ability to generalize
beyond memorized lexeme pairs.
\end{abstract}

\section{Introduction}
Modern natural language processing systems, particularly Large Language Models,
achieve strong performance across many natural language tasks, yet their
understanding of fine-grained lexical relations, the meaning of words and the
relations between them, remains limited; in this area, models often diverge
from human performance.

% across a wide range of tasks, especially with the development of Large
% Language Models (LLMs). By using language as input and output, these systems
% appear to possess strong human language understanding. This ability has been
% evaluated many times, finding limits to their understanding of the language
% they use. Lexical knowledge, the understanding of the words and of the
% relations between them, is a field of application where LLMs tend to diverge from human performance.

The Meaning-Text theory (MTT)
\cite{melcukVersLinguistiqueSensTexte1997} provides a linguistic framework that
decomposes language into multiple components, from measurable sound to abstract
meaning.
% : phonology, surface 
% morphology, deep morphology, surface syntax, deep syntax, and semantics. Going
% from phonology to semantics corresponds to language understanding, while the
% other way corresponds to language generation.
This theory allows for a universal
fine-grained description and analysis of multiple areas of linguistic
competency.
%, and is meant to be universal to all languages.

Within this framework, relations between lexical units are modeled using
lexical functions (LFs), which map a lexical unit to other units expressing a
semantic relation. For example, the \textit{Magn} function encodes
intensification, such as $Magn(\text{beautiful}) = \text{gorgeous, stunning}$.

% For a given function 
% $f$, for all lexical units $x$ with the associated meaning portrayed by $f$,
% $f(x)=y$ give the different possible values. For example, the $Magn$ function,
% which expresses the increase in intensity, can be applied to "beautiful",
% outputting $Magn(\text{beautiful}) = \text{gorgeous, stunning, amazing}$. These
% lexical functions can be used to create a fine-grained structured description
% of a language's lexicon.

Structured lexical representations such as LF offer a principled way to model
abstract lexical relations and may help language models move beyond
memorization of lexeme pairs. However, prior work has primarily used LFs for
evaluation rather than for improving models' lexical competence.

% might have the
% potential to help language models' capacities to better understand abstract
% lexical relations. More generally, structured lexical resources, such as
% ontologies describing the terminology of a specific domain like criminal law,
% have been used to support language understanding and generation task, for
% example with automatic translation.



This project aims to investigate whether language models already encode lexical
functions as abstract lexical relations, and if they are able to
generalize this abstraction or simply memorize lexeme pairs without
understanding the relations between them. To do so, we construct a dataset
based on lexical functions, and inject this knowledge into three pretrained
language models. We then evaluate whether this structured lexical training
improves lexical understanding and general linguistic competence.

% To do so, we will create a dataset inspired by previous work, using lexical
% functions as the primary input. We will then proceed to inject this knowledge
% into three pretrained language models and evaluate them on a 
% linguistic competence benchmark. More details about the task will be discussed
% in Section \ref{sec:task}.



\section{Related Works}\label{sec:relwork}
In order to evaluate language models' natural language understanding (NLU), the
enough generality to quantify models linguistic capabilities
\citep{niSurveyLargeLanguage2025}. GLUE \cite{niSurveyLargeLanguage2025}
proposed the first unified benchmark, giving a single score to quantify a
model's performance. Following in the steps of GLUE,
FLUE, a French specific benchmark composed of six linguistic evaluation tasks,
was proposed \citep{leFlauBERTUnsupervisedLanguage2020}. Inspired by this work,
but wanting to add more diversity and subtility to the evaluation, COLE was
created, proposing a 23 tasks evaluation benchmark with a single aggregated
score \citep{beaucheminCOLEComprehensiveBenchmark2025}. 

While general NLU benchmarks are useful, they do not allow for a fine-grained
evaluation of a specific linguistic evaluation, such as lexical capabilities.

The French Lexical Network (fr-LN) \citep{11403/lexical-system-fr/v3.2} is a
formal French lexicon following the MTT, describing a total of $29,976$ lexical
units. LFs are used to describe the relations between all the lexical units,
creating a graph data structure. Using this resource,
\citet{liuAssessingBERTsSensitivity2024} proposed 
an evaluation procedure to assess CamemBERT's \cite{martin-etal-2020-camembert}
understanding of French idioms. \citet{petrovQALFFrenchQuestionAnswering2025}
proposed the Q\&A-LF benchmark to evaluation models' lexical knowledge using
LFs. Using a fixed template prompt design for 15 selected lexical functions,
the authors found that tested language models generally excel when needed to use
morphological evidence, but struggle on tasks involving world conception and
creative reasoning, demonstrating the absence of structured semantic
mapping. Similarly, \citet{liEvaluationCompetenceLexicale2025} used 82 LFs to
compared language models.

These MTT specific evaluation procedures tend to remove rare LFs in their
testing, in order to have a more general score. To our knowledge, no prior work
proposed to enhance models with LFs knowledge.


\section{Task and Dataset} \label{sec:task}


Our primary task involves the extraction and structuring of lexical knowledge
to fine-tune a Large Language Model (LLM). We aim to teach the model specific
Lexical Functions derived from the Meaning-Text Theory.

As mentionned in the Section \ref{sec:relwork}, previous work on LFs and
language models created evaluation datasets, which are not suited for
optimization strategies. Based on previous work, we will create a structured
dataset using fr-LN \citep{11403/lexical-system-fr/v3.2} in JSON format. This
structure will include part of speech (POS), specific sense to handle polysemy,
relevant LFs, and an in-context example of the vocable. All of this data is
available in fr-LN; our work will be to structure it accordingly. As some LFs
have few instances in fr-LN, we plan to set a minimum occurrence threshold to
ensure consistency.

% To address the potential scarcity of data, we also plan to experiment with
% synthetic data generation techniques to augment the dataset. All synthetically
% generated entries will undergo manual validation to ensure linguistic
% correctness. We consider the development of these reliable augmentation
% strategies for structured lexical knowledge to be a contribution of our work.

\section{Methodology}
We will employ a parameter-efficient fine-tuning approach to verify if
structural injection improves competence.

We will be comparing three models : Mistral-7B
\citep{jiangMistral7B2023} for its modern and strong multilingual architecture;
Qwen2.5-14B-Instruct \citep{qwen2} for its instruction-tuned multilingual
model, and FlauBERT \citep{leFlauBERTUnsupervisedLanguage2020}, a strong
French-specific model. This diversity in models with manageable size for
academic experimentation will allow us to have a
broader understanding of language model's capacity regarding LFs.

Our primary objective is to compare the performance of the fine-tuned models
against the base models to see if explicit training yields better understanding
of lexical relations. Regarding the training techniques, we considered classic
full fine-tuning but will likely proceed with Low-Rank Adaptation (LoRA). LoRA
is significantly more cost-effective and computationally efficient, allowing us
to inject specific Meaning-Text Theory constraints without the prohibitive
resource costs of updating all model parameters. We will compare two main
configuration. Prompting the original Mistral-7B to perform tasks without
weight updates and the Mistral-7B model updated with our dataset via LoRA
\citep{LoRA}.

\section{Baselines and Evaluation}
Using the dataset as described in Section \ref{sec:task}, we will perform a two
step evaluation procedure.

First, by dividing the dataset in a $80\%, 20\%$ train-test split, we will
perform a supervised-learning inspired evaluation. Using
\citet{petrovQALFFrenchQuestionAnswering2025}'s evaluation procedure, we will
create fixed query template for every LF in our dataset. For example, a prompt
for the $S_0$ function could be "\textit{What noun corresponds to the verb
<<Verb>>?}". This strategy allows for an automatic evaluation procedure
on all terms in the test split. The prompts will also ask for a single lexical
unit answer.

Exact Match (EM) and Contain Match (CM)
\cite{petrovQALFFrenchQuestionAnswering2025} will be used to quantify the
model's answers; CM will be used to consider whether the model's answer
contains one of the desired outputs, since models tend to not respect the
"single lexical unit answer" condition. These metrics will require normalization of the
model's output, and will be set based, as there can be multiple good answers to
one LF input. Accuracy and F1 score will also be used to further analyze the
model's answers.

The second evaluation step will be using the COLE benchmark
\cite{beaucheminCOLEComprehensiveBenchmark2025} in order to assess whether
learning LFs allowed for a better general language understanding.


% Bibliography entries for the entire Anthology, followed by custom entries
%\bibliography{anthology,custom}
% Custom bibliography entries only
\bibliography{proposal}

\end{document}
